\chapter{State of the art}
\label{cap:contexto}

This chapter presents the scientific and technological framing of the work. It first characterises
the audience the system addresses and the limitations that define the problem; it then examines the
solutions currently available; and it finally surveys the technical areas from which the components
used are drawn, situating the contribution of this dissertation in relation to each of them. Whoever reads it will know that the retail investor's three questions already have answers in the literature, and why no free tool brings them together.

\section{The retail investor and the attention problem}
\label{sec:ctx_investidor}

This section characterises the documented behaviour of the audience the system addresses, since it
is the limitations of that behaviour, and not the size of the market, that delimit the problem
treated here.

Chapter~\ref{cap:introducao} established the scale of the universe in question. The
characterisation that is missing is behavioural, and it begins with an observation that runs
counter to the current portrayal of the retail investor as an agent who only reacts in panic.
During the fall of March 2020, retail investors on a commission-free platform increased their
aggregate positions rather than reducing exposure \autocite{welch2022robinhood}. This is an
audience that decides regularly, and that therefore benefits from better information at the moment
of the decision.

The behavioural finance literature establishes that those decisions depart systematically from the
rational ideal \autocite{barberis2003behavioral}. Two regularities bear directly on the design of
an alerting system. The first is loss aversion \autocite{kahneman1979prospect}, already stated in
the previous chapter, whose consequence for this work is that the cost of the absence of context is
asymmetric: it is greater precisely when the information is unfavourable.

The second is the attention effect, and it is the one that grounds the problem.
\textcite{barber2008glitters} establish that retail investors are net buyers of the stocks that
capture their attention, identified by three observable indicators, namely presence in the news,
abnormal trading volume and an extreme return on a single day. The mechanism the authors propose is
that of the search problem: faced with a very large number of alternatives, the investor does not
assess them all and the choice falls on the subset that has been made visible.

The authors delimit the reach of the claim, which concerns an aggregate imbalance between purchases
and sales, without extending to a universal behaviour, and report the conclusion that justifies this work: the
buying pattern so formed does not generate superior returns. The empirical basis is records from
United States brokerages between 1991 and 1999, prior to trading by mobile phone and to the absence
of commissions, so what is assumed transferable is the mechanism and not the measured values. The
centrality of attention in this process is such that it became measurable through the volume of
internet searches \autocite{da2011attention}.

The three indicators identified by the literature are all computed by this system, but the status
of each differs. News and the extreme movement are the two triggers described in
Chapter~\ref{cap:introducao} and can, on their own, give rise to an alert. Volume is computed and
presented as context, and does not interrupt the user. The reason for this asymmetry is empirical
and is reported in Section~\ref{sec:av_sintese}: combining volume with the remaining signals
improved triage in only one of three budgets considered, a result that does not support introducing
a third route of interruption. The idea of fusing signals of distinct origins into a single
indicator comes from a real-time information panel with that architecture
\autocite{worldmonitor2026}, suggested by the co-supervisor of this work, and enters the
dissertation in the condition in which it was evaluated, that is, as a measured and not adopted
alternative.

There is, on this point, field evidence that is more recent, closer to the object of this work, and
that points in an unfavourable direction. \textcite{liu2023alerts} followed 20\,904 retail
investors on a mutual fund platform in Taiwan, between January 2017 and June 2020, of whom 932
activated the platform's own price alerts. Estimating the effect by difference in differences, they
quantify it as a deterioration in portfolio performance of the order of one percentage point over
six months. The two mechanisms they identify are an increase of about 7\% in the number of
transactions, and the fact that the funds the investor adjusted after adoption performed worse than
those left untouched. The authors read this as the alert directing attention towards short-term
volatility and inducing excessive trading.

The result is uncomfortable for a dissertation that builds an alerting system, which is why it
appears in the body of the text and not in a footnote. It is worth being exact, even so, about what
was measured. The alert studied fires when the price of a fund crosses a threshold set by the
investor and reports the current price, with no indication of cause; it is the threshold alert that
Section~\ref{sec:ctx_ferramentas} encounters again in brokerage applications. The mechanism of harm
the authors describe, which is feedback trading on price rather than the use of information about
the value of the firm, is precisely what takes hold when the recipient is handed a number and
nothing else. The authors themselves conclude that improving access to information is not enough,
and that the capacity to process it must also be addressed.

This dissertation takes that result as a reason for the design rather than as validation of the
design it already had. Two decisions follow from it. The first is that the alert carries the
evidence that justifies it, instead of reporting only that a threshold was crossed. The second is
that the number of alerts is capped by explicit decision, and not by whatever the data produce, as
described in Section~\ref{sec:sis_politica}. Whether these two decisions suffice to avoid the
effect \textcite{liu2023alerts} measured remains to be demonstrated, and
Chapter~\ref{cap:conclusoes} takes it up in those terms.

On the side of financial institutions, the adoption of artificial intelligence has ceased to be
experimental, as documented by the survey of firms in the sector conducted by the
\textcite{ccaf2026aifs}. That survey does not measure the population of products aimed at the
retail investor, and therefore does not support claims about it; the characterisation of that
population rests on what the providers themselves declare, and is the object of the following
section.

\section{Financial information tools for retail}
\label{sec:ctx_ferramentas}

This section examines the solutions currently available, assessing each one against a fixed
criterion, namely the three questions stated in Chapter~\ref{cap:introducao}.
Table~\ref{tab:ctx_ferramentas} presents the result of that assessment.

\begin{table}[!htbp]
\caption[Existing tools against the three questions]{The categories of existing tool assessed
against the three questions of Chapter~\ref{cap:introducao}. The classification \emph{partial}
indicates that the tool supplies raw material for the answer and leaves the interpretation to the
user.}
\label{tab:ctx_ferramentas}
\centering
\small
\begin{tabular}{L{0.22\textwidth} C{0.11\textwidth} C{0.15\textwidth} C{0.12\textwidth}
                L{0.17\textwidth}}
\toprule
\tabhead{Tool} & \tabhead{Is it unusual?} & \tabhead{Company or market?} &
\tabhead{Has it happened?} & \tabhead{Cost} \\
\midrule
Brokerage application & no & no & no & included \\
News aggregator & no & no & no & free \\
Sentiment application & partial & no & no & free or low \\
Automated portfolio manager & no & no & no & percentage of the portfolio \\
Professional terminal & yes & yes & yes & not published \\
\midrule
\textbf{This work} & \textbf{yes} & \textbf{yes} & \textbf{yes} & \textbf{free} \\
\bottomrule
\end{tabular}
\end{table}

The table makes the essential point. The only existing tool that answers all three questions is
also the most expensive of all. The problem is therefore not scientific in nature, but one of access: the
answers are known and are conditional on a subscription incompatible with the investor profile in
question.

Three categories deserve individual examination, since each fails for a distinct reason and those
reasons inform decisions taken in the following chapters.

The brokerage application is the most immediate instrument of access to the price, and it is the
one whose public documentation describes the smallest set of indicators. It displays the day's
percentage change and a price chart. It does not indicate whether that change is large for that
particular stock, does not separate the company component from the market component, and has no
memory of comparable events.

The price alerts of these applications share the limitation in a more pronounced form. They fire
when a fixed threshold is crossed, a criterion that ignores the volatility proper to each stock and
therefore produces frequent alerts in agitated companies and rare ones in calm companies. They
further communicate that a level was crossed, and never why. This observation determines the
methodological option described in Section~\ref{sec:met_detecao}, which replaces the absolute
threshold with a measure relative to the recent norm of each company, and whose difference is
quantified in Section~\ref{sec:av_qi1}.

Sentiment applications classify news as positive or negative and present the result of that
classification. What they supply is a reading of the news item and not of the movement: a
strongly negative classification indicates that something has happened, and does not establish
whether the price change is large for that particular stock, which is the first question. It is
in this restricted sense that Table~\ref{tab:ctx_ferramentas} credits them with a partial
answer, and they go no further. Section
\ref{sec:av_qi2} additionally presents a measurement of our own according to which the topic of a
news item and the direction of the movement that follows it coincide less often than intuition
suggests, which limits the reach of a sentiment score as an indicator of impact.

Automated portfolio managers answer above all a different question, that of the allocation of
capital across a set of assets. The literature recognises genuine, though not uniform, benefits:
\textcite{dacunto2019robo} measure an improvement in diversification and a reduction in volatility
among investors who held fewer than five stocks before joining, whereas for those who already held
more than ten the effect on diversification is close to nil. The containment of behavioural biases
is the most consistently reported benefit \autocite{dacunto2019robo, cardillo2024robo}. The
difference with respect to this work is one of object and not of transparency: an automated manager
recommends a portfolio and executes it, whereas this system makes no recommendations and delivers
evidence about a concrete event. Moreover, providing investment advice is a regulated activity, a
boundary that the present work avoids by design, as discussed in Section~\ref{sec:met_etica}.

There are, finally, recent products that claim to answer exactly the question treated in this
dissertation. Robinhood Cortex, announced in March 2025, states on the provider's own page that it
is intended to answer the question of why a stock is going up or down on a given day
\autocite{robinhood2025cortex}. Google Finance's key moments, presented in June 2026, propose to
explain why a stock moved \autocite{google2026finance}. This is the same question, posed by
organisations with incomparably greater resources and user bases.

The difference lies in what the respective pages declare to be delivered. Both describe a summary
in natural language, and neither describes the presentation of individual past cases, with dates,
measured similarities and observed behaviour of the price after each of them, which is what this
work delivers. The distinction claimed is not one of fluency, but of verifiability, and it concerns
what the pages declare: the products were not tested, so nothing is asserted about what they
actually deliver.

Up to this point the section has examined products. The academic literature is missing, and leaving
it out would rest the claim of a gap on the inspection of vendor pages, which is a fragile basis.
The class of system at issue here, which is watching for events and notifying the retail investor, has
been treated in published research for some two decades, and the most developed example is the
\emph{MoFiN DSS} of \textcite{muntermann2009ubiquitous}. The parallel is close to the point of
being uncomfortable: it likewise addresses the retail investor, likewise starts from the
observation that this investor reacts to news more strongly and less quickly than the institutional
one, is likewise built and evaluated under the design science research paradigm that
Section~\ref{sec:met_metodologia} adopts, and likewise publishes its evaluation in a journal
article.

It is worth being precise about what that system does, because the positioning of this work follows
from that precision. Over 425 announcements by German firms published between 2003 and 2005, and
minute-by-minute intraday price series, the author fits two regression models: one estimates the
price effect that will follow the announcement, and the other estimates how long that effect will
last. The announcement is pushed to the investor when the estimated effect exceeds the average
effect observed. The evaluation is an \emph{ex ante} simulation of consumer surplus, that is, it
measures the economic opportunity the notification opens, and not the quality of what is
communicated. The user is told that an announcement is relevant and is not told why it is.

Three consequences for this dissertation. The first is that the gap cannot be stated as an absence
of alerting systems for the retail investor, because they exist and are published; relative to this
system, it lies in epistemic stance, between estimating what will happen and explaining what
already happened, and as a difference in what is delivered alongside the alert. The second is that
the relevance measure of \textcite{muntermann2009ubiquitous} is an abnormal return \emph{in
absolute value}, adjusted for the market and corrected for the usual behaviour of the stock itself,
that is, magnitude and not direction. It is in essence the same quantity that
Section~\ref{sec:met_detecao} measures, with two differences that matter: here it is measured after
the fact, and it serves to explain rather than to anticipate. That a work published in 2009 arrived
at magnitude, and not direction, as the operational definition of a relevant event for this
audience is independent support for a choice this dissertation reached by another route. The third
is that that work evaluated its artefact by simulation and explicitly deferred the study with real
users to future research, exactly the limitation Chapter~\ref{cap:conclusoes} states about this
one. This is no excuse, but it locates the problem: it is a difficulty of the field and not an
isolated omission.

Other works move towards explaining observed price movements. The \emph{Squawk Bot} of
\textcite{dang2020squawk} jointly learns price series and text to select news associated with
price behaviour. The \emph{DeepTrust} of \textcite{chan2022deeptrust} combines anomaly detection,
retrieval of Twitter posts and assessment of their reliability. The latter is a dissertation
made available on arXiv, whose evaluation concentrates on two episodes; its existence precludes
presenting the explanation of anomalies through information retrieval as a novelty of this work.

The \emph{Stock Pattern Assistant} of \textcite{neela2025spa}, also made available on arXiv,
segments prices into monotonic runs, associates events by temporal proximity and proposes a
narrative generated from those objects, without prediction or causal attribution. Its demonstration
uses events from a constructed sample, so it does not amount to continuous evaluation on captured
news. Explaining without predicting is therefore a shared choice. This dissertation is positioned
through the combination of the three questions, the evidence delivered and the evaluation of
components in operation, without claiming priority for that choice.

The most reasonable objection to the work consists in asking why a general-purpose artificial
intelligence assistant would not solve the same problem. No comparison with any such assistant was
carried out, so what follows is a design argument and not a result. An assistant without external
access to market data lacks three elements. The first is the recent volatility of the stock, which
allows one to determine whether a change is large. The second is the separation between the market
component and the company-specific component. The third is a base of past cases over which to
perform a retrieval.

On this last point, a language model recalls cases seen during training, which does not constitute
retrieval and is not verifiable by the reader. There are large models specialised in financial
text, whose capability is documented \autocite{li2023llmfinance}. The output of any of them remains
a statement that is read with ease and not checked.

The explainability literature gives reasons to distrust it. \textcite{lipton2018mythos} argues that
interpretability does not designate a single property and that a system should declare which of
them it offers. \textcite{rudin2019stop} argues that an explanation produced over an opaque model
is an approximation that may mislead. Neither deals with text generated by language models, so the
transposition is the author's and what follows from it is a design reason and not a result of the
literature. This is why language never produces facts in this system: the numbers come from the
components that computed them, and the text is composed from those, as described in
Section~\ref{sec:sis_explicacao}.

\section{Anomaly detection in financial time series}
\label{sec:ctx_anomalias}

This section situates the first decision of the system in the field of anomaly detection and states
the families of method considered.

Anomaly detection is an established area, with a consolidated taxonomy that separates methods by
the way they define normality \autocite{chandola2009anomaly, pang2021deep}. Three families are
relevant for financial time series.

The statistical family defines normality from the local distribution of the data themselves,
estimated in a rolling window, and flags as anomalous the observations that depart from it beyond a
threshold. The explanation of a detection coincides with the computation that produced it, which
makes it fully communicable to the user. In exchange, it assumes that the distribution of returns
is locally stable.

The unsupervised learned family defines normality from the geometric structure of the data.
Isolation Forest \autocite{liu2008isolation} identifies as anomalous the points that a random
partition isolates with few cuts; the Local Outlier Factor \autocite{breunig2000lof} identifies
them by their reduced density relative to that of their neighbours. Both dispense with labels and
capture multivariate patterns that a rule over a single series does not capture, at the cost of the
decision ceasing to be statable in one sentence.

The deep family learns a model of normality through a network \autocite{pang2021deep}. It offers
greater representational capacity, requires larger volumes of data, and departs from the
explainability requirement that structures this work.

A practical constraint conditions this whole choice. In neighbouring financial domains, such as
fraud detection, unsupervised methods occupy a central position precisely because labelled
anomalies are scarce and the patterns change over time \autocite{ahmed2016financial}. The same
condition holds here: there is no dataset in which it has been recorded, day by day, which
movements would justify an alert. This absence determines the way in which
Chapter~\ref{cap:avaliacao} can evaluate the component, and is the reason why the evaluation rests
on a criterion derived from the data themselves rather than on human judgement.

The assumption of local stability on which the statistical family rests is confronted and not
hidden. Conditional volatility models, introduced by \textcite{engle1982arch} and generalised by
\textcite{bollerslev1986garch}, formalise volatility clustering, that is, the tendency of calm and
agitated periods to occur in runs. Both papers formalise the phenomenon over inflation series, so
its transposition to stock returns is not assumed: Section~\ref{sec:av_qi1} measures it on the data
of this work, comparing the equally weighted window with a variant that assigns greater weight to
recent observations. A rolling window follows the phenomenon in part, since the norm against which
each day is compared is re-estimated daily from the near past. The choice does not eliminate the
problem, and the measurement referred to quantifies what it costs to keep it.

The comparison between the statistical family and the learned family does not remain at the level
of argument. Chapter~\ref{cap:avaliacao} runs both on the same data and reports the result
obtained.

\section{Text representation and semantic retrieval}
\label{sec:ctx_texto}

This section surveys the ways of representing text that make it possible to compare news by
meaning, and states the criterion for evaluating retrieval that follows from them.

Comparing news by meaning depends entirely on the representation adopted, and that representation
went through three generations, sketched in Figure~\ref{fig:ctx_geracoes} and detailed in
Table~\ref{tab:ctx_texto}.

\begin{figure}[htbp]
\centering
\begin{tikzpicture}[
  font=\small, >={Stealth[length=2.4mm]},
  g/.style={draw, rounded corners, align=center, text width=3.0cm, minimum height=3.3cm,
            inner sep=5pt, fill=black!3},
]
  \node[g] (a) {\textbf{Word counts}\\[3pt]
    \footnotesize each word is\\ \footnotesize an isolated symbol\\[3pt]
    \footnotesize \emph{``rose'' and ``gained''}\\ \footnotesize \emph{share no relation}};
  \node[g, right=0.85cm of a] (b) {\textbf{Word vectors}\\[3pt]
    \footnotesize words used in\\ \footnotesize similar ways lie close\\[3pt]
    \footnotesize \emph{one sense}\\ \footnotesize \emph{per word}};
  \node[g, right=0.85cm of b] (c) {\textbf{Sentence vectors}\\[3pt]
    \footnotesize meaning depends\\ \footnotesize on context\\[3pt]
    \footnotesize \emph{representation}\\ \footnotesize \emph{adopted in this work}};
  \draw[->, thick, black!45] (a) -- (b);
  \draw[->, thick, black!45] (b) -- (c);
\end{tikzpicture}
\caption[Three generations of text representation]{The three generations of text representation.
Each one resolves a limitation identified in the previous one, and the third is the one that allows
two sentences to be compared directly by the position of their vectors.}
\label{fig:ctx_geracoes}
\end{figure}

The first generation represents a document by the frequency of the terms it contains. The vector
space model \autocite{salton1975vsm} is transparent and computationally cheap, and it is blind to
synonyms: the terms ``rose'' and ``gained'' are, for this representation, symbols without relation.
Within the same generation there is a second line, coming from a distinct probabilistic framework,
that of ranking functions of which BM25 \autocite{robertson2009bm25} is the current
representative; both remain strong baselines in information retrieval and share that blindness. In
financial text there is also a specialised variant, the domain lexicons, which correct the fact
that generic sentiment dictionaries misread words that are common in financial language
\autocite{loughran2011liability}. They are entirely transparent, and are limited to the vocabulary
that was compiled.

The second generation represents each word as a vector in a space where proximity reflects similar
use, learned from local context \autocite{mikolov2013word2vec} or from global co-occurrence
\autocite{pennington2014glove}. This representation brings synonyms together, resolving the
previous limitation, and introduces a new one: each word receives a single vector, regardless of
the sentence in which it occurs.

The third generation removes that limitation. The Transformer \autocite{vaswani2017attention} made
it feasible to process the complete sentence with long-range dependencies, and \gls{BERT}
\autocite{devlin2019bert} produced contextual representations, in which the vector of a word
depends on those surrounding it. An additional step is needed for the problem treated here:
comparing two sentences with a model of that nature would require processing them jointly, once per
pair, which is not practicable over a base with tens of thousands of cases. \gls{SBERT}
\autocite{reimers2019sbert} resolves this by producing one vector per sentence, computed only once,
so that the comparison becomes an arithmetic operation between vectors that already exist.

The contribution of \gls{SBERT} that matters to this work does not lie in the architecture but in
the training objective: the vectors are learned so that the distance between two of them reflects
similarity of meaning. The authors report that an encoder of the \gls{BERT} family used directly to
compare sentences produces weak representations for that purpose, at times falling behind the
simple average of word vectors \autocite{reimers2019sbert}. It follows that an encoder tuned to
another task does not inherit this property by virtue of having been trained in the appropriate
domain.

\begin{table}[!htbp]
\caption[Text representation approaches]{The approaches considered for representing financial news,
with the representation each one produces and the implications of adopting it.}
\label{tab:ctx_texto}
\centering
\small
\begin{tabular}{L{0.22\textwidth} L{0.28\textwidth} L{0.36\textwidth}}
\toprule
\tabhead{Approach} & \tabhead{Representation} & \tabhead{Strengths and limitations} \\
\midrule
Term counts \newline \autocite{salton1975vsm, robertson2009bm25}
  & Weighted word frequency
  & Transparent and fast; blind to synonyms \\
\addlinespace
Financial lexicons \newline \autocite{loughran2011liability}
  & Word lists with assigned polarity
  & Fully transparent; limited to the compiled vocabulary \\
\addlinespace
Word vectors \newline \autocite{mikolov2013word2vec, pennington2014glove}
  & One fixed vector per word
  & Brings synonyms together; a single sense per word \\
\addlinespace
Sentence vectors \newline \autocite{devlin2019bert, reimers2019sbert}
  & One vector per sentence, sensitive to context
  & Compares sentences directly; larger model, less legible internally \\
\addlinespace
Domain models \newline \autocite{liu2020finbert, huang2023finbert}
  & As above, tuned to financial text
  & Domain vocabulary; superiority on this task is not automatic \\
\addlinespace
Large models \newline \autocite{li2023llmfinance}
  & Text generation over broad knowledge
  & High capability; expensive, opaque and outside the free-of-charge constraint \\
\bottomrule
\end{tabular}
\end{table}

For financial text there are models trained specifically for the domain
\autocite{liu2020finbert, huang2023finbert}. The natural assumption is that they are superior on
this task for that reason. The assumption is not automatically true, and
Section~\ref{sec:av_qi2} measures it rather than assuming it.

The work belongs to the wider literature of natural language processing applied to finance. The
review of \textcite{kearney2014textual} covers the state of that area up to 2014, dominated by
domain lexicons and by supervised classifiers over word counts; the contextual sentence
representations used in this work are later and do not appear in that review. There is also a
substantial line of research that uses text to anticipate price movements
\autocite{xing2018nlffsurvey, ding2015deep}. The present work stands outside that line by design
choice: text is used here to locate comparable earlier cases, whose outcome has already occurred
and has already been measured, and not to generate a forecast.

One recent work in that line deserves individual examination, being the closest to what is done
here. \textcite{du2024contrastive} train a triplet network over social media messages and price
series, and present to the user, alongside each prediction, the most similar historical case with
the same outcome and the most similar one with the opposite outcome. This is retrieval of analogues
used as explanation, over the same materials this work uses, and it is therefore the closest
published neighbour.

The differences are nonetheless fundamental, and worth stating because they delimit the
contribution. The label that organises the triplets is the direction of the next day's movement,
and the task evaluated is to predict it; magnitude enters only to define a dead band around zero.
The sentence encoder is used frozen, as the input layer to a convolutional network and a recurrent
network with attention, which are the ones actually trained, that is, the space in which news
items are compared is not changed. And the explanation is demonstrated through three case studies
with attention maps, with no measurement of its quality and no participants.

There is a result of their own in that work that deserves recording, because it concerns the
difficulty of the task and not the method. In the ablation study, applying contrastive learning to
the textual component alone degrades accuracy on all three benchmark datasets, falling below the
variant that uses no contrastive learning at all; the gains come from the quantitative component,
and the authors themselves conclude that textual information plays a complementary role. Anyone
seeking to use text to anticipate the direction of a movement has here, published, an indication
that the signal is weak; Section~\ref{sec:av_qi2} reaches a conclusion of the same sign by an
entirely independent route.

It is on those two differences that the fourth research question is built. If direction is the
weak signal that their ablation study shows, and if the comparison space there remains
unchanged, then two distinct questions are left open: whether the \emph{magnitude} of the
movement, which is not a direction and does not constitute a forecast, is learnable from
text, and whether the gain appears when the encoder itself is adjusted rather than a network
placed on top of it. Section~\ref{sec:met_qi4} describes how both are put to the test, with an
arm trained on magnitude and a replication arm trained on direction, and
Section~\ref{sec:av_qi4} reports what was obtained.

Representing each news item by a vector turns the identification of comparable cases into an
information retrieval problem. The proximity measure adopted is cosine similarity, whose mechanics
are described in Section~\ref{sec:met_recuperacao}. The operation is exact, since the whole case
base is scanned and the result ordered, so it does not depend on any approximation. At
substantially larger scales, exhaustive search can be replaced by an approximate index
\autocite{johnson2021faiss} without changing the representation, which is a path for growth and not
a present necessity.

Since retrieval returns an ordered list, its evaluation resorts to order-sensitive measures
\autocite{manning2008ir}. The one adopted in this work is precision at the first $k$ results, that
is, the fraction of those $k$ results that are relevant. The definition of relevance, which is
where an evaluation of this nature is decided, is established in Section~\ref{sec:met_recuperacao}
together with the baselines against which the result is compared.

\subsection{Retrieval with generation, and why the system stops before it}
\label{sec:ctx_rag}

The architecture in current use for answering a question from a body of documents is
retrieval-augmented generation, in which a retriever selects passages and a language model writes
the answer from them \autocite{lewis2020rag}. The literature on the family is extensive and has
been systematised \autocite{huang2026ragsurvey}.

The system described here executes the first half of that architecture and stops before the second,
and the reason is the one that runs through the work. Retrieval returns verifiable objects: a
headline, a date, a measured similarity and an observed price movement. Generation converts those
objects into a statement in prose, whose faithfulness to what was retrieved is itself a property to
be evaluated, and which the recipient of this system has no means to confront. The option does not
follow from a technical limitation, and the distinction should be drawn clearly: adding generation
is possible and cheap; what is not cheap is the evaluation that it does not displace the meaning of
the evidence. No comparison with a complete architecture was carried out in this work, so what is
asserted here is a design reason and not a result. Section~\ref{sec:sis_ciclo} describes the
grounded generation layer that was in fact built and evaluated, and the reason why it is not
exposed.

\subsection{Financial news corpora and sector labels}

Evaluating retrieval over financial news requires a body of text aligned with prices, and that
alignment is what distinguishes a usable corpus from a collection of headlines. The dataset used in
this work is \gls{FNSPID} \autocite{dong2024fnspid}, which brings together news and the
corresponding price series and whose construction is documented.

Relevance between two news items is approximated, in this work, by belonging to the same sector.
The option replaces a non-existent human annotation with an objective criterion, and carries the
corresponding limitation. The validity of a grouping is measurable, by the separation between
groups \autocite{rousseeuw1987silhouettes} or by the agreement between alternative groupings
corrected for chance \autocite{vinh2010ami}. Neither of those measures was applied to the sector
map adopted. What the label guarantees is objectivity and reproducibility, and not that the sector
is the most informative partition of the corpus.

\section{Event studies and return decomposition}
\label{sec:ctx_evento}

This section presents the instrument with which the effect of an event on the price is measured, as
well as the statistical correction on which the separation between the company component and the
market component depends.

The instrument adopted is the event study, whose formulation goes back to
\textcite{fama1969adjustment}. The behaviour of the method over daily returns was examined by
\textcite{brown1985daily}, who compare alternative ways of estimating the expected return and
conclude that the current procedures, based on estimating the market model by least squares, remain
well specified in short windows around the event. The reach of this conclusion is delimited by the
authors themselves: the simplest variant, which subtracts only the historical mean of the stock,
loses statistical power and may become misspecified when the event dates coincide. It is this
result that motivates the order of magnitude of the one, three and five day windows adopted in this
work. The transposition is not direct, and the difference is declared in
Section~\ref{sec:met_recuperacao}, according to which the value presented to the user is the raw
return and not the abnormal return that those authors study.

The existence of a statistical relation between published text and market movements is a long
standing observation \autocite{tetlock2007media}. The appropriate verb is that of evidence of a
statistical relation, and not that of a demonstration of causality.

One methodological point conditions the second question of the work, that of the separation between
the company and the market. That separation requires estimating the sensitivity of each stock to
the market, and the estimate is noisy, especially in volatile companies or those with a short
history. \textcite{blume1971risk} establishes that the coefficients so estimated tend to regress
towards the mean between successive periods, from which followed the practice of shrinking the
estimate with a fixed weight. This work adopts instead the shrinkage weighted by the imprecision
of the estimate itself \autocite{vasicek1973beta}, whose formulation and effect are presented
in Section~\ref{sec:met_decomposicao}.

The theoretical framing that justifies the refusal to forecast also belongs to this area. The
efficient market hypothesis, in its semi-strong form, holds that prices already reflect publicly
available information \autocite{fama1970efficient}, from which it follows that future movements
should not be systematically predictable from that information. It is one of the reasons why the
system measures what has already happened rather than projecting what is to come; the second, of a
practical nature, is that only the past is verifiable by the reader.

\section{Case-based reasoning}
\label{sec:ctx_casos}

This section frames the third technique of the system within an established line of artificial
intelligence and delimits the part of that line the work implements.

Case-based reasoning \autocite{aamodt1994cbr} approaches a new problem through the retrieval of
similar earlier problems and the use of what is known about them, dispensing with exclusive
recourse to general domain knowledge. The source organises the method into a cycle of four
processes, namely retrieve, reuse, revise and retain, and considers all four necessary to the
complete resolution of a problem.

Figure~\ref{fig:ctx_cbr} shows the cycle and marks where the system stops.

\begin{figure}[htbp]
\centering
\begin{tikzpicture}[font=\footnotesize,
  cheio/.style={draw=black!75, line width=0.9pt, rounded corners=2pt, fill=black!6,
            minimum width=2.45cm, minimum height=1.05cm, align=center},
  vazio/.style={draw=black!45, densely dashed, line width=0.8pt, rounded corners=2pt,
            fill=white, minimum width=2.45cm, minimum height=1.05cm, align=center,
            text=black!55},
  leg/.style={font=\scriptsize, text=black!65, align=center},
  seta/.style={-{Stealth[length=4pt]}, draw=black!50},
]
  \node[cheio] (n0) at (0.0,0) {\textbf{retrieve}\\[1pt]\scriptsize finds similar\\ past cases};
  \node[cheio] (n1) at (3.15,0) {\textbf{reuse}\\[1pt]\scriptsize draws on what\\ was measured in them};
  \node[vazio] (n2) at (6.3,0) {\textbf{revise}\\[1pt]\scriptsize adapts the outcome\\ to the new case};
  \node[vazio] (n3) at (9.45,0) {\textbf{retain}\\[1pt]\scriptsize stores the solution\\ as a future case};
  \draw[seta] (n0) -- (n1);
  \draw[seta] (n1) -- (n2);
  \draw[seta] (n2) -- (n3);
  \draw[seta, densely dashed, draw=black!35] (n3) -- ++(0,-0.95) -- (0,-0.95) -- (n0);
  \node[leg] at (1.58,0.95) {the system implements};
  \node[leg] at (7.87,0.95) {the system stops here};
  \draw[draw=black!30] (4.72,1.15) -- (4.72,-0.62);
  \node[leg, text width=8.2cm] at (4.72,-1.7) {because adapting would produce an expectation about what comes next, which is the forecast the work refuses};
\end{tikzpicture}
\caption[The case-based reasoning cycle, and where the system stops]{The four processes that \textcite{aamodt1994cbr} consider necessary to the complete resolution of a problem. The system runs the first two in full and stops before the next two. The interruption is not a shortfall of implementation: \emph{revise} adapts the outcome of a past case to the present one, and that adaptation is an expectation about the future, that is, exactly the forecast the founding constraint refuses. What the system presents is the measured outcome of each case, as it occurred. The first two processes, run in full and with the decision left to the person, correspond to what \textcite{kolodner1991aiding} calls case-based decision aiding.}
\label{fig:ctx_cbr}
\end{figure}
The technique described in Section~\ref{sec:met_recuperacao} implements the first two and stops
deliberately before the rest: it retrieves similar cases and presents what was measured on them,
without adapting that outcome to the present case. Adaptation would produce an expectation about
what is going to happen, which is precisely the forecast the work refuses. The partial
implementation of an established cycle is, in this case, a consequence of the founding constraint
and not a shortcoming.

This stopping point has, moreover, a name of its own in the literature, and recognising it changes
its status. \textcite{kolodner1991aiding} describes \emph{case-based decision aiding} as a
methodology in which the system augments the person's memory by supplying analogous cases, and in
which it is the person, not the machine, who decides, using those cases as guidance. The starting
point is an observation from psychology: people reason well by analogy but remember the right cases
poorly, whereas computers remember well. Put that way, what this system does ceases to be the
execution of half a cycle and becomes the full execution of a distinct methodology, proposed in
1991 by the person who defined the area. The difference is not one of pride: a partial
implementation explains what it lacks, whereas a methodology declares what it delivers and to whom
the decision belongs.

It is worth recording, in the opposite direction, what case-based reasoning has actually done in
this domain, so that the originality claimed is no larger than it is.
\textcite{oh2007cbrmonitoring} build a daily indicator of the condition of the Korean financial
market that classifies each day as a stable, unstable or crisis period, retrieving cases from a
base built on the 1997 crisis. This is case-based reasoning applied to finance, and it is
instructive for two reasons. The first is the object: the unit classified is the whole market, in a
judgement about systemic risk, and not the movement of one company following a news item. The
second is the destination of the retrieved cases, which serve to produce a label automatically; the
article does not describe their presentation to anyone. Showing the cases, which is what is done
with them here, is precisely what is left undone there.

The option finds additional support on the side of the reader. Research on explanation in the
social sciences establishes that people look for contrastive and selective reasons, and not for an
exhaustive enumeration of everything that contributed to an outcome
\autocite{miller2019explanation}. A small set of concrete cases, with dates and outcomes,
corresponds better to that expectation than a coefficient. It is for that reason that each alert
presents a few precedents and the exact quantities that support the decision, rather than the
totality of the values computed.

There is a second requirement, on the same side, that conditions not the content of the explanation
but the frequency with which it is delivered. A recipient exposed to excessive warnings stops
processing them, an effect measured on retail investors \autocite{liu2023alerts} and discussed
in Section~\ref{sec:ctx_investidor}. The consequence is that the decision policy forms part of the
design. An alert that is not read is worth the same as an alert that was not sent,
and that is why Chapter~\ref{cap:sistema} treats the suppression criteria with the care it devotes
to the methods.

\section{Explainability and trust in automated systems}
\label{sec:ctx_xai}

This section situates the explainability requirement that runs through the work and presents the
cautions that the literature associates with the presentation of explanations.

The demand that a system show how it reached a conclusion is an area in its own right, with
consolidated reviews \autocite{arrieta2020xai, adadi2018peeking}. The area divides into two routes,
shown in Figure~\ref{fig:ctx_xai}: transparent models, interpretable in themselves, and
\emph{post-hoc} methods, which explain an opaque model already trained and which can be organised
by the type of object they explain \autocite{guidotti2018survey}.

\begin{figure}[htbp]
\centering
\begin{tikzpicture}[
  font=\small, >={Stealth[length=2.4mm]},
  r/.style={draw, rounded corners, align=center, text width=3.5cm, minimum height=1.1cm,
            inner sep=4pt, fill=black!3},
  t/.style={draw, rounded corners, align=center, text width=3.5cm, minimum height=1.4cm,
            inner sep=4pt, fill=black!8, very thick},
  o/.style={draw, rounded corners, align=center, text width=3.5cm, minimum height=1.4cm,
            inner sep=4pt, fill=black!3},
]
  \node[r] (raiz) {\textbf{Explaining a decision}};
  \node[t, below left=1.1cm and 0.35cm of raiz] (tr) {\textbf{Transparent model}\\
    \footnotesize the explanation is the computation};
  \node[o, below right=1.1cm and 0.35cm of raiz] (ph) {\textbf{Post-hoc method}\\
    \footnotesize explains an opaque model};

  \node[below=0.4cm of tr, align=center, font=\footnotesize, text width=3.6cm]
    (trx) {rule, linear regression,\\ count over the series};
  \node[below=0.4cm of ph, align=center, font=\footnotesize, text width=3.6cm]
    (phx) {\gls{LIME}, \gls{SHAP}};

  \draw[->, thick, black!45] (raiz) -- (tr);
  \draw[->, thick, black!45] (raiz) -- (ph);
  \draw[-, black!35] (tr) -- (trx);
  \draw[-, black!35] (ph) -- (phx);
\end{tikzpicture}
\caption[The two routes to explainability]{The two routes distinguished by the literature
\autocite{arrieta2020xai, guidotti2018survey}: building models that are transparent by design, or
explaining an opaque model after training. This work follows the first.}
\label{fig:ctx_xai}
\end{figure}

The two most widely used \emph{post-hoc} methods operate in distinct ways. \gls{LIME} fits a simple
model in the neighbourhood of an individual prediction, locally approximating the behaviour of the
opaque model \autocite{ribeiro2016lime}. \gls{SHAP} attributes a prediction to the variables that
produced it using Shapley values \autocite{lundberg2017shap}. Both are generic and both are useful.

Both share, however, a property that is decisive in this context. Faced with decisions of serious
consequence, \textcite{rudin2019stop} argues for abandoning opaque models in favour of models
interpretable by construction, on the grounds that every \emph{post-hoc} explanation is itself an
approximation, and an approximation can lead to mistaken readings. The user receives in that case
two objects whose relation they are not in a position to assess, namely a decision and an
explanation that may not correspond to it.

This work follows the conservative route. Instead of explaining an opaque model \emph{post hoc}, it
resorts to components in which the explanation and the computation coincide. A departure from a
recent norm, an additive split whose sum reproduces the observed movement, and a set of earlier
cases with the outcomes measured are objects presentable in full. The explanation is not a second
object produced over the decision: it is the decision, stated at length.

Two cautions from the literature close this section, and both condition the organisation of
Chapter~\ref{cap:avaliacao}.

The first is that interpretability does not designate a single property
\autocite{lipton2018mythos}, so a system should declare the type of interpretability it offers
rather than claiming the term without qualification. This work offers transparency of mechanism,
that is, the user can follow the computation that led to the decision, and does not offer a causal
explanation of the event in the world.

The second concerns the effect of the explanation itself on the person who receives it. Trust in
automated systems fails in two symmetric directions, with the user ignoring a correct warning or
accepting an incorrect one, so the objective is not to maximise trust but to make it appropriate to
the real capability of the system \autocite{lee2004trust}. \textcite{bansal2021whole} add that the
presence of an explanation increases the probability of accepting the system's answer whether it is
right or wrong, that is, that the observed gain does not come from the person becoming better at
distinguishing the two cases. The design consequences that follow are described in
Section~\ref{sec:met_etica}.

These cautions are general. On the concrete choice made by this dissertation, which is to explain by
cases, there exists empirical evidence gathered with people that points the other way, and
relegating it to a footnote would be the cheapest way of neutralising it.

\textcite{waa2021xai} compared rule-based explanations, example-based explanations and no
explanation, in two experiments with forty-five participants each, in a decision aid for insulin
dosing. Rule-based explanations significantly improved identification of the decisive factor;
example-based explanations were indistinguishable from no explanation at all, neither in factor
identification ($p = 0.796$) nor in self-reported understanding ($p = 0.283$). The explanation the
authors give for the result is the part that matters most here: both styles, they write, provide
only the details relevant to a single decision, and not the underlying rationale or causality.

Within the financial domain itself, and with non-experts, there are two recent studies.
\textcite{cau2023logicstyle} compare three explanation styles in a stock trading task and conclude
that the abductive and deductive styles improve performance when the model is confident, by
comparison with the inductive style, which is the example-based style and therefore the one closest
to what this system delivers. \textcite{bertrand2023featurebased} test an automated life-insurance
advisor with feature-based explanations and conclude that these improve neither understanding nor
appropriate reliance, by comparison with giving no explanation at all, and that the more
conversational format increases the user's trust, sometimes to the user's detriment.

Three consequences follow for this dissertation. The first is that the choice of case-based
explanation, justified in Section~\ref{sec:ctx_casos} by the literature on explanation in the
social sciences \autocite{miller2019explanation}, has no support in the measurements with people
that exist, and in one case has evidence from the domain itself against it. The second is that the
diagnosis of \textcite{waa2021xai} names exactly what this system adds to the cases: an additive
decomposition whose sum reproduces the observed movement, which is the underlying computation whose
absence the authors hold responsible for the result. The system does not deliver cases alone; it
delivers cases and the arithmetic that accompanies them. Whether that is sufficient is an open
empirical question, and the work does not answer it. The third, and the most consequential, is that
these results make the study with users more necessary and not less: there is published literature
pointing against the founding hypothesis, and it is in those terms that
Section~\ref{sec:con_limitacoes} states its absence.

The second study supports, by contrast, a decision already taken. If the conversational format is
what produces excessive trust, then not exposing the conversational layer described in
Section~\ref{sec:sis_inteligencia} ceases to be merely engineering prudence and acquires empirical
support within the domain.

Furthermore, claims of interpretability demand rigorous evaluation and are not established by
declaration \autocite{doshivelez2018considerations}. It is for that reason that
Chapter~\ref{cap:avaliacao} measures what is measurable by automatic procedure and explicitly
identifies the part that remains unevaluated for lack of a controlled study with users.

\section{Machine learning in a production environment}
\label{sec:ctx_producao}

This section brings together the literature relating to the learned component of the system,
organised according to the three questions its deployment raises.

The first question is that of determining against what stronger reference a simple model should be
compared. Ensembles of trees trained by gradient boosting \autocite{friedman2001gbm} capture
interactions and non-linear effects that a logistic regression does not capture, and therefore play
that role in Chapter~\ref{cap:avaliacao}. It is not claimed that they constitute the best existing
method, since that would be a claim about an entire field.

The second question concerns the honesty of the value presented. A score is only interpretable as a
probability if it is calibrated, and the raw scores of many classifiers are not; \emph{post-hoc}
calibration over a held-out validation block corrects that distortion
\autocite{niculescu2005calibration}. The same source does not, however, favour the option taken in
this work: it identifies logistic regression as already well calibrated to begin with and concludes
that, above roughly a thousand calibration cases, isotonic regression matches or exceeds the
sigmoid. The deployed model is a logistic regression with a validation block far above that
threshold, so both statements apply to it. The option is justified by measurement and not by
authority: the comparison of the two variants over the same validation set is reported in
Table~\ref{tab:av_alternativas}, on page~\pageref{tab:av_alternativas}.

An alternative route was considered and not adopted. Conformal prediction returns a set of labels
with a distribution-free guarantee, valid in finite samples \autocite{vovk2005algorithmic,
angelopoulos2023conformal}. Two clarifications determine that it is not applicable under the
conditions of this work. The coverage of the current variant is marginal: it holds on average over
cases, and not case by case. And the guarantee requires exchangeability between the calibration
block and the test block, an assumption that the chronological split with embargo of
Section~\ref{sec:met_triagem} refuses by construction. Section~\ref{sec:con_futuro} returns to the
method as a path to explore.

The third question is that of the persistence of performance over time. A deployed model faces
change in the distribution of the data it observes, and the concept drift literature establishes
the taxonomy of that change and the strategies for detecting and adapting to it
\autocite{gama2014survey}. This work measures drift on its own data rather than assuming it absent.

There is, finally, a cost that is captured by no quality metric. \textcite{sculley2015debt}
establish that machine learning components accumulate hidden maintenance, in particular through
entanglement between parts and through unstable, undeclared data dependencies.
Chapter~\ref{cap:conclusoes} returns to this point with an example measured in the present system,
in which part of the learning cycle stopped without any error being recorded.

\section{Synthesis and positioning of the work}
\label{sec:ctx_sintese}

This section locates the work with respect to the review presented and states the scope decisions
that follow from it.

The review allows the space in which the work is inscribed to be delimited precisely, and that
space is not a technical void. Every component used exists, is published and is evaluated, as
Table~\ref{tab:ctx_posicionamento} makes explicit.

\begin{table}[!htbp]
\caption[Positioning of the work with respect to the literature]{Every component of the system
exists in the literature in isolation. What does not exist is their combination under the
zero-cost constraint, with the evidence delivered to the user at the moment the statement is
presented.}
\label{tab:ctx_posicionamento}
\centering
\small
\begin{tabular}{L{0.24\textwidth} L{0.30\textwidth} L{0.34\textwidth}}
\toprule
\tabhead{Question} & \tabhead{What the literature offers} & \tabhead{What the target audience
lacks} \\
\midrule
Is it unusual?
  & Anomaly detection, statistical and learned \autocite{chandola2009anomaly, liu2008isolation}
  & A measure relative to the company's own norm, communicated in accessible language \\
\addlinespace
Is it the company or the market?
  & Event study and sensitivity estimation \autocite{brown1985daily, vasicek1973beta}
  & Availability outside a professional terminal \\
\addlinespace
Has it happened before?
  & Sentence representation and retrieval \autocite{reimers2019sbert, manning2008ir}
  & A case base with measured outcomes, and not recalled ones \\
\addlinespace
Why believe it?
  & Explainability and the trust literature \autocite{rudin2019stop, lee2004trust}; measurements
    with people that do not confirm the gain \autocite{waa2021xai, cau2023logicstyle,
    bertrand2023featurebased}
  & The evidence delivered alongside the statement and verifiable at the moment of reading;
    still to be demonstrated with users \\
\bottomrule
\end{tabular}
\end{table}

The gap is therefore not one of method, but of integration under constraint. The products that
answer the three questions are professional terminals, unavailable free of charge or at a price
accessible to the retail investor. Their price is not published in a citable form, and it is
unavailability, and not a concrete value, that supports the argument. The free products that claim
to answer the third deliver a statement without the evidence that supports it.
The academic literature includes integrated systems and proposals for explaining past movements,
as Section~\ref{sec:ctx_ferramentas} documents. The distinction from MoFiN DSS includes the absence
of prediction; with DeepTrust and Stock Pattern Assistant, that choice is shared. The contribution
must therefore be assessed through the combination of rarity, decomposition and precedents, the
operating constraints and the evaluation of the evidence delivered. The reviewed works delimit
this contribution; they do not demonstrate a universal absence of comparable systems.

This observation determines the distinction that runs through the work and that
Chapter~\ref{cap:sistema} takes up when justifying the absence of a language model in the
composition of the messages. A generated sentence without links to facts presents itself to the reader
in a fluent register and gives no means of confronting it with what supports it. A sentence
composed from computed objects is a statement accompanied by the evidence, in which every value
presented points back to the measurement that produced it and the reader can confirm each element
at the moment of reading. This is the property the system seeks to preserve, and it is what fixes
the nature of the contribution. It is not the formulation of a new method. It is the selection,
integration and critical evaluation of existing methods in a working system, under declared
constraints, and with the results that do not favour the options taken included in the evaluation.

Four scope decisions follow from this positioning, and each one deliberately excludes a class of
functionality.

\begin{description}
  \item[Answer the three questions, and only those.] No functionality is incorporated that does not
        serve one of them, which keeps the system at a size compatible with explaining it in full.
  \item[Produce no price forecasts.] Neither direction, nor target, nor recommendation. Besides the
        verifiability reason stated in Chapter~\ref{cap:introducao}, the theoretical framing of
        Section~\ref{sec:ctx_evento} applies \autocite{fama1970efficient}.
  \item[Manage no portfolios and give no advice.] The system does not know the user's positions and
        does not tell them what action to take, which avoids collecting personal data and keeps the
        work outside the boundary of regulated financial advice.
  \item[Use free sources exclusively.] The constraint makes the system accessible to the audience
        it addresses and turns the resulting limitations into measured and reported quantities.
\end{description}

The last decision has a real cost that should be stated without mitigation. News arrives late and
the coverage of events is incomplete. Chapter~\ref{cap:sistema} quantifies both limitations and
Chapter~\ref{cap:conclusoes} assesses their consequences.
